%% kapa: exam=1 round=27 year=2016 period=2 form=A subject=appraisal_law subject_ko=감정평가관계법규 seq=30 number=30
%% source: Q-Net 감정평가사 기출문제 (cst00309) artlSeq=5169402 file=extracted/1008/5169402/제27회 감정평가사 1차 문제/제27회 감정평가사 2교시문제(A형).pdf
%% conversion: kordoc 4.24.51 -> tools/gen1.py (auto); review=proofread-sonnet; flags=none
\begin{question}{30}
공간정보의 구축 및 관리 등에 관한 법령상 지상 경계의 결정기준을 옳게 연결한 것을 모두 고른 것은?
\begin{framedtext}
ㄱ. 연접되는 토지 간에 높낮이 차이가 없는 경우: 그 구조물 등의 중앙 \par
ㄴ. 토지가 해면 또는 수면에 접하는 경우: 최대만조위 또는 최대만수위가 되는 선 \par
ㄷ. 도로ㆍ구거 등의 토지에 절토된 부분이 있는 경우: 경사면의 중앙 \par
ㄹ. 공유수면매립지의 토지 중 제방 등을 토지에 편입하여 등록하는 경우: 바깥쪽 하단부 \par
ㅁ. 연접되는 토지 간에 높낮이 차이가 있는 경우: 그 구조물 등의 하단부 \par
\end{framedtext}
\begin{choices}
  \item ㄱ, ㄹ
  \item ㄷ, ㅁ
  \item ㄱ, ㄴ, ㄹ
  \item ㄱ, ㄴ, ㅁ
  \item ㄴ, ㄷ, ㄹ, ㅁ
\end{choices}
\end{question}
\answer{4}
